\documentclass[11pt]{article}

\usepackage[a4paper,margin=30mm]{geometry}
\usepackage{amsmath,amssymb,amsthm,mathtools}
\usepackage[T1]{fontenc}
\usepackage{lmodern}
\usepackage{microtype}
\usepackage[hidelinks]{hyperref}

\newtheorem{theorem}{Theorem}[section]
\newtheorem{lemma}[theorem]{Lemma}
\newtheorem{remark}[theorem]{Remark}
\newcommand{\D}{\mathbb D}
\newcommand{\N}{\mathbb N}
\newcommand{\R}{\mathbb R}
\newcommand{\qb}{q_b}
\newcommand{\Rb}{\mathcal R_b}

\hypersetup{
  pdftitle={A third magnetic Neumann level below the Dirichlet ground state on the disk},
  pdfauthor={Xiaojian Zeng},
  pdfsubject={An exact Dirichlet--Neumann eigenvalue comparison for the magnetic unit disk},
  pdfkeywords={magnetic Laplacian, unit disk, Dirichlet--Neumann inequality, min--max principle}
}

\title{A third magnetic Neumann level below the Dirichlet ground state on the disk}
\author{Xiaojian Zeng\\
\small PBC School of Finance, Tsinghua University, Beijing 100190, China\\
\small \texttt{xiaojianzeng@foxmail.com}}
\date{}

\begin{document}
\maketitle

\begin{abstract}
Let $H_b=(-i\nabla-b(-y,x)/2)^2$ be the spinless magnetic Laplacian on the
unit disk, and let $\lambda_k(b)$ and $\mu_k(b)$ denote its Dirichlet and
magnetic-Neumann eigenvalues, respectively, counted with multiplicity.  We
prove that $\mu_3(b)\leq\lambda_1(b)$ for every $b>0$.  The previously
unresolved part of this comparison is the bounded-field interval
$0<b\leq4$.  There, the positive radial function
$(1-r^2)e^{-r^2/2}$ gives a piecewise explicit lower bound on the
Dirichlet ground energy.  We compare that bound with the Rayleigh quotients
of three mutually orthogonal magnetic-Neumann form-domain states.  Two
states work throughout the interval, while the third is chosen from one of
the two angular-momentum states meeting the Dirichlet ground state at
$b=2$.  The differences factor into elementary expressions whose signs are
fixed on $0<b\leq2$ and $2\leq b\leq4$.  Min--max then avoids any spectral
branch ordering or truncation.  Combining this variational certificate with
the published magnetic-Pauli counting argument for $b>4$ proves the
all-field statement.
\end{abstract}

\noindent\textbf{Keywords:} magnetic Laplacian; unit disk; Dirichlet--Neumann
inequality; eigenvalue comparison; min--max principle

\medskip
\noindent\textbf{Mathematics Subject Classification (2020):}
35P15; 35J10; 81Q10

\section{Introduction}

Inequalities between Dirichlet and Neumann spectra measure how much the
low-energy state count changes when a hard boundary is replaced by a
reflecting one.  For the homogeneous magnetic Laplacian on a bounded convex
planar domain, Lotoreichik proved the index-shifted comparison
$\mu_{k+1}\leq\lambda_k$ \cite{Lotoreichik2025}.  On the disk, numerical
evidence in the same work suggested the sharper first-index relation
$\mu_3(b)\leq\lambda_1(b)$.  The cited paper proved this relation for
$b>4$ and identified an exact three-branch crossing at $b=2$, but left the
continuous interval $0<b\leq4$ at the numerical level.  The nonmagnetic
analogue and its geometric extensions form a related but distinct problem;
see, for example, \cite{Rohleder2025}.  Recent detailed work on the magnetic
Neumann disk has focused on its lowest eigenvalue and angular-momentum
crossings \cite{HelfferLena2025}.

The purpose of this note is to close the disk comparison by an exact
low-field argument.  Its ingredients are deliberately elementary: a
ground-state representation for a positive Dirichlet comparison function,
four explicit trial functions, angular orthogonality, and the min--max
principle.  In particular, the proof never has to determine which radial
branch is the third eigenvalue.

Let
\[
  \D=\{(x,y)\in\R^2:x^2+y^2<1\},\qquad
  A_b(x,y)=\frac b2(-y,x),\qquad b>0,
\]
and set $H_b=(-i\nabla-A_b)^2$.  The Dirichlet and magnetic-Neumann
realizations are defined by the common quadratic form
\[
  \qb[u]=\int_{\D}|(-i\nabla-A_b)u|^2\,dx\,dy
\]
on $H_0^1(\D)$ and $H^1(\D)$, respectively.  Their eigenvalues, repeated
according to multiplicity, are denoted by $\lambda_k(b)$ and $\mu_k(b)$.
Our main result is the following.

\begin{theorem}\label{thm:main}
For every $b>0$,
\[
  \mu_3(b)\leq\lambda_1(b).
\]
\end{theorem}

The normalization is dimensionless.  For a disk of radius $R$, dilation
replaces $b$ by $bR^2$ and multiplies both spectra by $R^{-2}$.

\section{Angular decomposition and a Dirichlet lower bound}

With the convention $u(r,\theta)=f(r)e^{im\theta}$, $m\in\mathbb Z$, the
radial expression in $L^2((0,1);r\,dr)$ is
\begin{equation}\label{eq:fiber}
  T_{b,m}=-\frac{d^2}{dr^2}-\frac1r\frac d{dr}
           +\left(\frac mr-\frac{br}{2}\right)^2.
\end{equation}
The regular solution behaves as $O(r^{|m|})$ at the origin.  Dirichlet data
give $f(1)=0$.  Since $A_b$ is tangent to the circle, magnetic-Neumann data
reduce in this gauge to $f'(1)=0$.  Distinct angular sectors are orthogonal
both for the $L^2$ inner product and for $\qb$.

We use the established disk fact that the Dirichlet ground energy is the
ground energy of the $m=0$ fiber \cite[Sec.~5]{Lotoreichik2025}.

\begin{lemma}\label{lem:dirichlet}
For $0<b\leq4$,
\begin{equation}\label{eq:L}
  \lambda_1(b)\geq L(b):=
  \begin{cases}
    5+b^2/4,&0<b\leq2,\\
    6,&2\leq b\leq4.
  \end{cases}
\end{equation}
\end{lemma}

\begin{proof}
Put
\[
  \phi(r)=(1-r^2)e^{-r^2/2}.
\]
This function belongs to the radial Dirichlet form domain and is positive
for $0<r<1$.  Direct differentiation in \eqref{eq:fiber} gives
\begin{equation}\label{eq:pointwise}
  \frac{T_{b,0}\phi}{\phi}
  =6+\frac{b^2-4}{4}r^2.
\end{equation}
For a smooth radial form-core function $f$, integration by parts yields the
ground-state representation
\begin{equation}\label{eq:gst}
\begin{aligned}
 &\int_0^1\left(|f'|^2+\frac{b^2r^2}{4}|f|^2\right)r\,dr
 -\int_0^1\frac{T_{b,0}\phi}{\phi}|f|^2r\,dr \\
 &\hspace{35mm}=
 \int_0^1\phi^2\left|\left(\frac f\phi\right)'\right|^2r\,dr\geq0.
\end{aligned}
\end{equation}
The inequality extends to the radial Dirichlet form domain by closure.  For
$0<b\leq2$, the infimum of the right-hand side of
\eqref{eq:pointwise} is $5+b^2/4$; for $2\leq b\leq4$, it is $6$.
The radial characterization of the Dirichlet ground state now gives
\eqref{eq:L}.
\end{proof}

\section{An explicit three-dimensional trial space}

For a nonzero $u\in H^1(\D)$, write
\[
  \Rb[u]=\frac{\qb[u]}{\|u\|_{L^2(\D)}^2}.
\]
Two trial functions will be used for the whole interval $0<b\leq4$:
\[
  u_0(r,\theta)=1,\qquad
  u_1(r,\theta)=re^{i\theta}.
\]
Direct radial integration gives
\begin{equation}\label{eq:twoR}
  \Rb[u_0]=\frac{b^2}{8},\qquad
  \Rb[u_1]=4-b+\frac{b^2}{6}.
\end{equation}
Consequently $\Rb[u_0]\leq2$ and $\Rb[u_1]\leq4$ on this interval,
so both are strictly below $L(b)$.

The third state changes at $b=2$.  Define
\begin{align*}
  v_-(r,\theta)&=re^{-r^2/2}e^{-i\theta},\\
  v_+(r,\theta)&=r^2(3-r^2)e^{-r^2/2}e^{2i\theta}.
\end{align*}
Both are regular $H^1(\D)$ functions.  At $b=2$, their radial parts satisfy
$T_{2,-1}v_-=6v_-$ and $T_{2,2}v_+=6v_+$, respectively, as well as the
Neumann condition at $r=1$ \cite[Sec.~5]{Lotoreichik2025}.  Set
\begin{equation}\label{eq:rho}
  \rho_\pm=
  \frac{\int_0^1r^2|v_\pm(r)|^2r\,dr}
       {\int_0^1|v_\pm(r)|^2r\,dr}.
\end{equation}
Only the exact support bound
\begin{equation}\label{eq:rho-range}
  0<\rho_\pm<1
\end{equation}
will be needed.

\begin{lemma}\label{lem:third}
The trial state $v_-$ satisfies $\Rb[v_-]\leq L(b)$ for $0<b\leq2$,
and $v_+$ satisfies $\Rb[v_+]\leq L(b)$ for $2\leq b\leq4$.
\end{lemma}

\begin{proof}
Changing $b$ in \eqref{eq:fiber} changes only the potential
\[
  \left(\frac mr-\frac{br}{2}\right)^2
  =\frac{m^2}{r^2}-mb+\frac{b^2r^2}{4}.
\]
Using the exact eigenvalue $6$ at $b=2$ therefore gives
\begin{equation}\label{eq:vminusR}
  \Rb[v_-]=6+(b-2)+\frac{b^2-4}{4}\rho_-.
\end{equation}
For $0<b\leq2$, subtraction of the first branch of \eqref{eq:L}
factorizes as
\begin{equation}\label{eq:vminusfactor}
  \Rb[v_-]-\left(5+\frac{b^2}{4}\right)
  =\frac{b-2}{4}
   \left((\rho_--1)b+2(1+\rho_-)\right).
\end{equation}
The expression in parentheses is positive: it decreases with $b$, and at
$b=2$ it equals $4\rho_->0$.  Thus \eqref{eq:vminusfactor} is nonpositive.

For the $m=2$ state, the same potential shift gives
\begin{equation}\label{eq:vplusfactor}
  \Rb[v_+]-6=(b-2)
  \left(-2+\frac{b+2}{4}\rho_+\right).
\end{equation}
When $2\leq b\leq4$, \eqref{eq:rho-range} implies
\[
  -2+\frac{b+2}{4}\rho_+
  <-2+\frac{b+2}{4}\leq-\frac12.
\]
Hence \eqref{eq:vplusfactor} is also nonpositive.
\end{proof}

\begin{proof}[Proof of Theorem~\ref{thm:main} for $0<b\leq4$]
For $0<b\leq2$, take
\[
  \mathcal V_b=\operatorname{span}\{u_0,u_1,v_-\},
\]
whose elements have angular momenta $0,1,-1$.  For $2\leq b\leq4$,
take instead
\[
  \mathcal V_b=\operatorname{span}\{u_0,u_1,v_+\},
\]
with angular momenta $0,1,2$.  In either case $\dim\mathcal V_b=3$ and
angular orthogonality diagonalizes both the norm and the quadratic form.
The Rayleigh quotient of every nonzero element of $\mathcal V_b$ is thus a
weighted average of the three individual quotients.  Equations
\eqref{eq:twoR}, \eqref{eq:L}, and Lemma~\ref{lem:third} show that its
maximum is at most $\lambda_1(b)$.  The magnetic-Neumann min--max principle
then yields
\[
  \mu_3(b)\leq
  \max_{0\ne u\in\mathcal V_b}\Rb[u]
  \leq\lambda_1(b).
\]
The switch point $b=2$ and the endpoint $b=4$ are included.
\end{proof}

\section{The high-field interval}

For completeness, we recall how the published high-field result joins the
preceding argument.  The normalized flux of the constant field through the
unit disk is
\[
  \Phi=\frac{1}{2\pi}\int_{\D}b\,dx\,dy=\frac b2.
\]
The magnetic-Pauli counting theorem of Fournais, Frank, Goffeng, Kachmar,
and Sundqvist \cite[Prop.~2.1]{FournaisEtAl2025}, specialized to the disk in
\cite[Sec.~5]{Lotoreichik2025}, implies that $b>4$ gives at least three
magnetic-Neumann eigenvalues strictly below $b$.  Thus
\begin{equation}\label{eq:highN}
  \mu_3(b)<b\qquad(b>4).
\end{equation}

On the Dirichlet form domain, if
$\Pi_x=-i\partial_x+by/2$ and $\Pi_y=-i\partial_y-bx/2$, then
$[\Pi_x,\Pi_y]=ib$ and integration by parts gives
\[
  \qb[u]-b\|u\|^2
  =\|(\Pi_x+i\Pi_y)u\|^2\geq0.
\]
Consequently $\lambda_1(b)\geq b$.  Together with
\eqref{eq:highN}, this proves
$\mu_3(b)<\lambda_1(b)$ for $b>4$.  Combining it with the low-field proof
completes Theorem~\ref{thm:main}.

\section{Interpretation and limitations}

The theorem is a spectral statement for two self-adjoint boundary
realizations of the same spinless constant-field Hamiltonian.  It says that
replacing the hard-wall boundary by magnetic reflection puts at least three
levels, counted with multiplicity, at or below the hard-wall ground energy.
The proof is nonperturbative and uses no finite angular-momentum cutoff,
radial discretization, or numerical ordering of crossing branches.

The circular geometry is essential to the present certificate: it supplies
the angular-sector orthogonality and the explicit states at $b=2$.  No
conclusion is asserted here for noncircular domains, variable fields, Robin
conditions, Pauli operators with a Zeeman term, or comparisons involving
higher indices.

\section*{Statements and declarations}

\paragraph{Data and code availability.}
The result is analytic and uses no empirical data.  The accompanying source
package contains a symbolic verification script and its recorded output for
the differentiations, Rayleigh quotients, and factorizations used above.

\paragraph{Competing interests and funding.}
No competing-interest or funding declaration has been supplied for this
pre-submission draft.  The corresponding author must provide and approve the
applicable declarations before submission.

\paragraph{Artificial-intelligence assistance.}
OpenAI Codex was used in September 2026 as an assistive tool for literature
triage, proof search, algebraic checking, adversarial review, and manuscript
drafting.  It is not an author and cannot assume responsibility.  The named
author has not yet completed the independent scientific and editorial review
required before submission and must verify the argument, citations,
disclosures, and final text before assuming accountability for them.

\bibliographystyle{amsplain}
\bibliography{references}

\end{document}
